\chapter{绪论}

\section{研究背景与意义}

药物研发是一个周期很长、成本很高、风险也很大的过程,从最初的靶点发现到最后的临床试验,整个流程下来往往需要花费十几年甚至更长的时间,投入的资金也是以十亿美元来计算的\cite{di2019innovation}。在这个漫长的过程中,先导化合物的发现和优化是特别关键的一个环节,它直接决定了后续研发工作的成败。传统的药物分子设计主要依靠的是药物化学家的专业知识和经验积累,他们会根据对疾病机制的理解以及对受体结构的分析,来设计和合成具有潜在活性的化合物。不过这种方法存在一个很明显的问题,就是效率比较低,因为化学空间实在是太大了,据估计可能存在的类药分子数量达到了$10^{60}$这个级别\cite{bohacek1996art}，而人类能够实际合成和测试的分子数量只是其中很小的一部分。

为了解决这个问题,计算化学和计算机辅助药物设计(CADD)技术应运而生。这些技术通过建立数学模型和算法,能够在计算机上对大量的分子进行虚拟筛选和性质预测,从而缩小实验验证的范围,提高研发效率\cite{lyu2019ultra}。比如说定量构效关系(QSAR)方法,它就是通过统计分析已知的活性分子数据,建立起分子结构与生物活性之间的定量关系模型,然后用这个模型来预测新分子的活性\cite{muratov2020qsar}。还有分子对接技术,它可以模拟小分子配体与生物大分子受体之间的相互作用,评估结合的自由能,从而判断分子的潜在活性\cite{meng2011molecular}。这些计算方法在一定程度上缓解了传统实验方法的局限性,但是它们也面临着一些挑战。一方面,现有的预测模型的准确性还是有限的,特别是对于那些结构新颖的分子,预测结果往往不够可靠；另一方面,这些方法大多是基于局部搜索的策略,容易陷入局部最优解,难以在全局的化学空间中找到真正优秀的候选分子。

最近几年,人工智能技术的快速发展为药物研发带来了新的机遇。深度学习模型在图像识别、自然语言处理等领域取得了突破性的进展,也开始被应用到药物发现的各个环节中\cite{chen2018rise}。特别是在分子生成这个方向上,出现了很多基于深度学习的生成模型,比如变分自编码器(VAE)、生成对抗网络(GAN)、以及基于强化学习的方法等\cite{olivecrona2017molecular}。这些模型能够从已有的分子数据中学习化学空间的分布规律,然后生成具有特定性质的新分子结构。相比于传统的基于规则的方法,深度学习模型具有更强的表达能力和泛化能力,能够探索到更加广阔和复杂的化学空间。

在这些人工智能技术中,大语言模型(Large Language Model, LLM)的发展尤其引人注目。从GPT-3到GPT-4,再到各种开源的大模型如LLaMA\cite{touvron2023llama}、Qwen等,这些模型展现出了惊人的语言理解和生成能力\cite{achiam2023gpt}。它们不仅能够完成文本生成、机器翻译、问答系统等传统的自然语言处理任务,还能够在代码生成、逻辑推理、甚至是科学发现等方面发挥作用。大语言模型之所以能够做到这些,是因为它们在训练过程中接触到了海量的文本数据,从中学习到了丰富的世界知识和复杂的模式规律。这种强大的知识表示和推理能力,给很多领域都带来了启发,包括我们关注的药物分子设计领域。

把大语言模型应用到分子设计上,一个很自然的想法就是把分子的SMILES表示当成一种特殊的"语言"来处理。SMILES(Simplified Molecular Input Line Entry System)是一种用ASCII字符串来表示分子结构的记号系统,它能够把二维的分子图结构编码成一维的字符串,这样就能够利用自然语言处理的技术来进行处理和分析了\cite{weininger1988smiles}。事实上,已经有一些研究工作尝试用Transformer架构来学习SMILES序列的模式,并用于分子生成和性质预测\cite{schwaller2019molecular}。不过这些工作大多是专门训练一个针对分子的模型,而没有充分利用到预训练大语言模型已经在大规模通用语料上学到的知识和能力。

本研究就是要探索如何把现成的大语言模型直接应用到分子优化任务中,而不是重新训练一个专门的模型。我们的核心思路是用大语言模型的文本生成能力来替代传统进化算法中的变异操作,通过自然语言的提示词来引导模型理解优化的目标和要求,然后生成具有改进性质的新分子。这种方法有几个明显的优势：一是可以利用到大语言模型已经学到的丰富化学知识,不需要从头开始训练；二是通过自然语言的提示,可以灵活地调整优化的目标和约束条件,比传统的固定算子更加灵活；三是大语言模型本身具有很强的创造性,能够生成一些人类化学家可能想不到的新颖结构,有助于打破思维定式,发现新的化学空间区域。

从理论意义上来说,本研究探索了大语言模型在科学计算和工程设计领域的应用潜力,展示了自然语言处理技术与计算化学交叉融合的可能性。它为如何利用预训练模型的知识和能力来解决专业领域的复杂问题提供了一个具体的案例,对于推动人工智能技术在科学研究中的应用具有积极的意义。从实际应用的角度来看,如果能够证明这种方法是有效的,那么它就可以作为一个辅助工具,帮助药物化学家更快地找到有潜力的候选分子,缩短药物研发的周期,降低研发的成本。虽然目前这个方法还处于初步探索的阶段,距离真正的工业应用还有一定的距离,但是它所体现的思路和方向是值得深入研究和发展的。

\section{国内外研究现状}

在大语言模型应用于分子设计这个方向上,国内外的研究者都已经开展了一些相关的探索工作。从国外的研究来看,Molecular Transformer是最早引起关注的工作之一。Schwaller等人提出用Transformer模型来预测化学反应的产物,他们把反应物和试剂的SMILES作为输入,产物的SMILES作为输出,通过大量的反应数据进行训练,使得模型能够学习到化学反应的规律\cite{schwaller2019molecular}。这个模型在多个基准测试数据集上都取得了很好的效果,证明了Transformer架构在处理化学序列数据方面的有效性。在此基础上,后续的研究者又提出了各种改进的版本,比如加入注意力机制的解释性分析、引入反应条件的信息、以及扩展到逆合成分析等任务。

除了反应预测,直接用大语言模型来生成分子也是一个活跃的研究方向。ChemBERTa是一个专门为化学领域设计的预训练语言模型,它在数百万个分子SMILES数据上进行了掩码语言模型的预训练,学习到了分子的语法和语义特征\cite{chithrananda2020chemberta}。研究表明,经过预训练的ChemBERTa在下游的分子性质预测任务上,只需要少量的标注数据就能达到很好的性能,这体现了预训练-微调范式在化学领域的适用性。类似的工作还有MoLFormer,它采用了不同的预训练策略和模型架构,也在多个基准测试中表现优异\cite{ross2022molformer}。这些工作说明,针对化学领域进行专门的预训练,确实能够提升模型在相关任务上的表现。

不过上述这些工作都是要重新训练或者微调一个专门的模型,这需要大量的计算资源和标注数据。与之不同的是,有一些研究者开始探索如何直接使用通用的大语言模型来完成化学任务。比如Galactica模型,它是Meta公司推出的一个面向科学知识的大型语言模型,在训练数据中包含了大量的科学文献、教科书、以及数据库信息\cite{taylor2022galactica}。研究发现,Galactica能够在不进行专门微调的情况下,完成分子性质预测、反应预测、甚至是论文写作等任务,展现出了很强的零样本和小样本学习能力。这说明通用的大语言模型确实已经学到了一定的化学知识,可以直接应用于一些化学相关的任务。

在国内,相关的研究起步稍微晚一些,但是发展速度很快。清华大学、北京大学、中国科学院等机构的研究团队都在这个方向上开展了工作。比如清华大学的团队提出了一种结合知识图谱和大语言模型的分子生成方法,通过知识图谱提供结构化的化学知识,大语言模型负责生成和推理,两者互补,提高了生成分子的质量和多样性\cite{wang2024knowledge}。北京大学的团队则专注于多目标分子优化,他们设计了复杂的提示词策略,让大语言模型能够同时考虑多个药物属性指标,生成综合性能更好的分子\cite{liu2024multi}。这些工作都体现了国内研究者在这个前沿领域的积极探索和创新精神。

在进化算法与机器学习结合方面,也有很多相关的研究。传统的进化算法包括遗传算法、进化策略、粒子群优化等,它们在组合优化、参数调优等问题上有着广泛的应用\cite{holland1992genetic,eiben2015introduction}。近年来,研究者开始尝试用机器学习模型来增强进化算法的各个组件,比如用神经网络来预测个体的适应度,减少昂贵的评估次数；用强化学习来选择变异算子,提高搜索的效率；用生成模型来初始化种群,提供更好的起点等\cite{karafotias2014parameter}。这些工作被称为"神经进化"或者"学习型进化算法",它们的核心思想是利用机器学习模型的泛化能力和预测能力,来弥补传统进化算法盲目搜索的不足。

本研究提出的方法可以看作是神经进化的一个具体实例,不过我们用大语言模型替代了传统的变异算子,这是一个比较新颖的尝试。与之前的工作相比,我们的方法有几个特点：一是完全依赖于预训练的大语言模型,不需要额外的训练或微调；二是通过自然语言提示来实现灵活的优化目标控制；三是保留了进化算法的种群机制和选择压力,确保搜索的方向性和有效性。这些特点使得我们的方法在实现简单性和应用灵活性方面具有一定的优势。

\section{研究内容与目标}

基于以上的研究背景和现状分析,本研究的主要目标是探索和验证大语言模型在分子进化优化中的应用可行性。具体来说,我们要回答以下几个核心问题：第一,大语言模型是否能够理解分子的SMILES表示,并根据提示生成化学上有效且具有改进性质的新分子；第二,不同的提示词策略对生成结果有什么样的影响,哪种策略更有效；第三,把大语言模型集成到进化算法框架中,能否实现有效的分子优化,其性能 compared to 传统方法如何；第四,大语言模型生成的随机性对优化过程有什么影响,如何控制和利用这种随机性。

围绕这些问题,本研究的具体内容包括以下几个方面：

一是构建一个完整的LLM驱动的分子进化优化系统。这个系统需要整合多个功能模块,包括分子表示和验证模块,用来处理和验证SMILES字符串的化学有效性；属性预测模块,用来计算分子的各种药物属性指标,如QED、LogP、分子量等；大语言模型交互接口,用来调用本地的Ollama服务,发送提示并接收生成的结果；提示词模板设计,根据不同的优化策略设计合适的提示词格式和内容；种群管理模块,负责维护分子种群,执行选择、更新等操作；以及结果分析和可视化模块,用来记录实验数据和生成图表。

二是设计和实现多种提示词策略。我们计划设计三种不同层次的提示策略,从简单到复杂,逐步增加提示的信息量和指导性。第一种是基础优化提示,只包含当前分子的SMILES、目标属性名称、以及当前的分数,要求模型生成改进的分子。第二种是带示例的少样本学习提示,在基础提示的基础上,增加几个成功优化的例子,让模型能够通过类比学习来理解优化的策略。第三种是鼓励多样性的探索提示,除了基本的优化要求外,还会告诉模型近期已经生成过的分子,要求它避免重复,尝试新的化学结构。通过对比这三种策略的效果,我们可以了解提示信息对模型生成行为的影响规律。

三是进行系统的实验验证和性能评估。我们会选择QED(定量估计药物相似性)作为主要的优化目标,因为它是一个综合性的指标,能够反映分子的药物开发潜力。实验会分为几个部分：首先是对比三种提示策略的优化效果,看哪种策略能够更快地提升QED分数；其次是分析种群的多样性变化,检查是否存在过早收敛的问题；然后是测试大语言模型生成的随机性,通过调整温度参数和控制多次运行,来评估结果的稳定性和可重复性；最后是与基线方法进行对比,虽然我们没有实现传统的进化算法作为对照,但是可以通过文献中的数据来进行定性的比较分析。

四是总结和提炼研究的发现和启示。在实验的基础上,我们会分析大语言模型在分子优化任务中的优势和局限性,探讨可能的改进方向和应用前景。同时也会反思研究过程中遇到的问题和挑战,为后续的研究工作提供参考和建议。

\section{论文组织结构}

本论文的其余部分组织如下：第二章介绍相关的理论基础和技术背景,包括分子表示方法、药物属性预测、进化算法的基本原理、以及大语言模型的工作机制等,为后续的方法设计和实验分析奠定理论基础。第三章详细阐述我们提出的LLM驱动的分子进化优化方法,包括系统的整体架构、各个模块的设计和实现、提示词策略的设计思路、以及关键的算法流程等。第四章呈现实验的结果和分析,包括实验设置、对比实验的结果、多样性分析、随机性测试、以及结果讨论等。第五章总结全文的主要工作和贡献,指出研究的不足之处,并对未来的研究方向进行展望。
